\documentclass[10pt,a4paper]{amsart}
\usepackage[margin=19mm]{geometry}
\usepackage{amsmath,amssymb,mathtools}
\usepackage[scr=rsfs,cal=euler]{mathalfa}
\usepackage{xfrac}
\usepackage[unicode,hidelinks,bookmarks=false]{hyperref}
\newcommand{\ie}{i.e.\ }
\newcommand{\cf}{cf.\ }
\newcommand{\resp}{respectively\ }
\newcommand{\Subsection}{\subsection}
\providecommand{\nobreakdash}{\nobreak\hbox{-}}
\setlength{\emergencystretch}{2em}
\multlinegap0pt
\usepackage{bm}
\usepackage{bbm}
\usepackage{dsfont}
\usepackage{xspace}
\usepackage{cancel}
\usepackage{mathtools}
\usepackage{cleveref}
\usepackage[shortlabels]{enumitem}
\usepackage{amsthm}
\makeatletter
\@ifundefined{theorem}{\newtheorem{theorem}{Theorem}}{}
\@ifundefined{lemma}{\newtheorem{lemma}[theorem]{Lemma}}{}
\makeatother
\newcommand{\NN}{\mathbb{N}}
\newcommand{\fm}{\mathfrak{m}}
\newcommand{\fn}{\mathfrak{n}}
\newcommand{\fq}{\mathfrak{q}}
\title[Bertini's theorem for $F$-rational $F$-pure singularities]{Bertini's theorem for $F$-rational $F$-pure singularities}
\author{Alessandro De Stefani, Thomas Polstra, Austyn Simpson}
\date{Source: arXiv:2509.04433v2 (CC BY 4.0)}
\begin{document}
\maketitle

\noindent\emph{Source and licence.} Bertini's theorem for $F$-rational $F$-pure singularities. Version arXiv:2509.04433v2 (CC BY 4.0), distributed under the Creative Commons Attribution 4.0 International licence, as recorded in the arXiv licence metadata for this exact version and shown on the abstract page https://arxiv.org/abs/2509.04433v2. Full rights notice: licenses/arxiv-2509.04433-CC-BY-4.0.txt.\par\smallskip

\noindent\textit{Editorial scope.} Counted unit: the Lemma whose source label is \texttt{B1:F-rational-F-pure} in the article (source0.tex lines 514--515, source label \texttt{B1:F-rational-F-pure}), delivered together with the article's own complete original proof of it (source0.tex lines 517--547, one \texttt{proof} environment of 31 lines). No mathematics is written, repaired or re-derived: statement and proof are the article's own text line for line. Required context retained uncounted: thm:F-sing-base-change (lines 396--402); thm:F-pure-base-change (lines 398--401). Every definition, display and label of those lines is retained unchanged. Omitted from this package and not redistributed: 1--395, 402--513, 516--516, 548--862. Nothing inside a retained excerpt is omitted or altered: every retained body is delivered exactly as the article's lines are written, so no omission receipt of any kind accompanies this package. Citations: the retained excerpts carry no citation command at all, so no bibliography entry of the article is transcribed and no reference is redistributed. Reference adaptation: none was needed and none was made; every reference inside the delivered unit resolves inside it, so no unresolved reference or citation is produced. Printed slips kept literal: no printed word, symbol, hypothesis or number of the counted statement or of its proof was altered. Layout and class: the article's own class is \texttt{amsart} with packages including geometry, inputenc, mathptmx, mathalfa, mathrsfs, setspace, amssymb, amsopn, amsmath, array, pdfsync, url, amsfonts, float; the wrapper is the lane's pinned \texttt{amsart} editorial wrapper at 10pt on A4, which restates the theorem-like environments the retained text needs and adds the article's own macro definitions that the retained text needs (\texttt{\textbackslash NN}, \texttt{\textbackslash fm}, \texttt{\textbackslash fn}, \texttt{\textbackslash fq}), with extra wrapper packages (bm, bbm, dsfont, xspace, cancel, mathtools, cleveref, enumitem). The delivered numbering is the unit's own. Assets: no retained span contains a figure environment, a graphics-inclusion command, a PDF-page inclusion, a table or a TikZ picture; no image file is copied into this package, the package declares no asset dependency and no image file is redistributed with it. Licence: version arXiv:2509.04433v2 of this article is distributed under the Creative Commons Attribution 4.0 International licence, as recorded in the arXiv licence metadata for this exact version and shown on the abstract page https://arxiv.org/abs/2509.04433v2. A full rights notice is delivered at licenses/arxiv-2509.04433-CC-BY-4.0.txt. Characters: every retained excerpt is pure ASCII, so the delivered engine, XeLaTeX with its default Latin Modern text font, reports no missing character.
\begin{theorem}\label{thm:F-sing-base-change}
    Let $(R,\fm)\to (S,\fn)$ be a flat local ring homomorphism between excellent local rings of prime characteristic $p>0$. Then:
\begin{enumerate}[label=(\roman*)]
    \item If $R$ is $F$-pure and $S/\fm S$ is a regular local ring (or more generally, Gorenstein and $F$-pure), then $S$ is $F$-pure;\label{thm:F-pure-base-change}
    \item If $R$ is $F$-rational and $S/\fm S$ is geometrically regular over $R/\fm$, then $S$ is $F$-rational. \label{thm:F-rat-base-change}
\end{enumerate}
\end{theorem}

\begin{lemma} \label{B1:F-rational-F-pure} Let $K \subseteq L$ be a finite extension of fields of characteristic $p>0$, and let $R$ be a $K$-algebra of finite type. Assume that the map $R \to S:=R \otimes_K L$ has regular fibers, and that $R$ is $F$-rational and $F$-pure. Then $S$ is $F$-rational and $F$-pure.
\end{lemma}

\begin{proof}
We have that $S$ is $F$-pure by \cref{thm:F-sing-base-change} \ref{thm:F-pure-base-change}. We now show that $S$ is also $F$-rational.

    Note that the extension $R \subseteq S$ is finite and free, and both $S$ and $R$ are Cohen--Macaulay domains. Let $\fn$ be a maximal ideal of $S$, so that $\fm:=\fn \cap R$ is a maximal ideal of $R$. Note that $\fn$ is a minimal prime of $\fm S$, and set $(A,\fm):=R_\fm \to S_\fn=:(B,\fn)$. Let $d$ be the common Krull dimension of the local rings $A$ and $B$.
    
    In order to show $F$-rationality of $B$, first assume that the extension $K \subseteq L$ is finite separable. We claim that, in this case, the map $A \to B$ is faithfully flat with geometrically regular closed fiber. Since both rings are excellent, this will imply that $B$ is $F$-rational, as desired, see \cref{thm:F-sing-base-change}. It is clear that the map is faithfully flat.
    Let $\ell$ be a finite field extension of $\kappa(\fm):=R/\fm$. Since $K \subseteq \kappa(\fm)$ is also finite, we have that $\ell$ is a finite extension of $K$. Note that
    \[
    S/\fm S \otimes_{\kappa(\fm)} \ell \cong (L \otimes_K \kappa(\fm)) \otimes_{\kappa(\fm)} \ell \cong L \otimes_K \ell.
    \]
    Since $K \subseteq L$ is finite separable, the ring $L \otimes_K \ell$ is a product of fields. This shows that the map $A \to S_W$, with $W=R \smallsetminus \fm$, has geometrically regular fibers, and therefore so does $A \to B$.

    Now assume instead that the finite field extension $K \subseteq L$ is purely inseparable. In particular, there exists $e_{0}\in\NN$ such that $L^{p^{e_0}} \subseteq K$. Then $S^{p^{e_0}} \subseteq R$, and 
    we also get that $B^{p^{e_0}} \subseteq A$.
    
    Let $\fq = (x_1,\dots, x_d)\subseteq A$ be a common system of parameters for $A$ and $B$. The $F$-rationality of $B$ will follow from the equality $\fq B =(\fq B)^*$. To that end, let $y\in (\fq B)^*$. There exists $c\in B^{\circ}$ so that
    \begin{align*}
        cy^{p^e} \in (x_1^{p^e},\dots, x_d^{p^e})B
    \end{align*}
    for all $e\gg 0$. Since $c^{p^{e_0}}$ and $y^{p^{e_0}}\in A$, we have by flatness of $A\to B$,
    \begin{align*}
        c^{p^{e_0}} y^{p^{e+e_0}}\in (x_1^{p^{e+e_0}},\dots, x_d^{p^{e+e_0}})B \cap A = (x_1^{p^{e_0}},\ldots,x_d^{p^{e_0}})^{[p^e]}A
    \end{align*}
    for all $e\gg 0$. It follows that
    \begin{align*}
        y^{p^{e_0}}\in \left((x_1^{p^{e_0}},\dots, x_d^{p^{e_0}})A\right)^* = (x_1^{p^{e_0}},\dots, x_d^{p^{e_0}})A\subseteq (x_1^{p^{e_0}},\dots, x_d^{p^{e_0}})B
    \end{align*}
    by flatness and the assumption that $A$ is $F$-rational. Since $B$ is $F$-pure, it follows that $y\in \fq B$, and thus $B$ is $F$-rational as claimed. 
    
    To handle the general case, write the finite extension $K \subseteq L$ as $K \subseteq L_1 \subseteq L$, where the $K \subseteq L_1$ is separable and $L_1 \subseteq L$ is purely inseparable. Since $S = R \otimes_K L \cong (R \otimes_K L_1) \otimes_{L_1} L$ we conclude that $S$ is $F$-rational by combining the two cases discussed above.
\end{proof}

\end{document}
